\documentclass{ximera}

\input{../../preamble.tex}

\title{Compose}

\begin{document}

\begin{abstract}
replacement
\end{abstract}
\maketitle





If we write a composition in function notation, $(f \circ g)(x) = f(g(x))$, then we see that the algebra of composition is replacement.





\begin{example} Composition


\begin{itemize}
\item Let $K(u) = \frac{u}{u-1}$ with its natural or implied domain: $(-\infty, 1) \cup (1, \infty)$. 

\item Let $T(w) = w^2 + 5w + 7$ with its natural or implied domain: \textbf{$\mathbb{R}$}.
\end{itemize}


Form the composition $K \circ T$.

\[        (K \circ T)(w) =     \frac{w^2 + 5w + 7}{(w^2 + 5w + 7)-1}   =    \frac{w^2 + 5w + 7}{w^2 + 5w + 6}  \]


$K$ can accept any number, except $1$.  Therefore we need to find out when $T = 1$ and take out the domain numbers where $1$ occurs.



\begin{align*}
T(w) & = 1   \\
w^2 + 5w + 7 & = 1 \\
w^2 + 5w + 6 & = 0   \\
(w+2)\left( \answer{w+3} \right) & = 0
\end{align*}


We need to remove $-2$ and $\answer{-3}$ from the domain of $K \circ T$, because $T(-2)=1$ and $T(-3)=1$ and $1$ cannot be an input into $K$.


The domain for $K \circ T$ is $(-\infty, -3) \cup (-3, -2) \cup (-2, \infty)$




\[        (K \circ T)(w)  =  K(T(w))  = \frac{w^2 + 5w + 7}{w^2 + 5w + 6}  =    \frac{w^2 + 5w + 7}{(w+2)(w+3)} \]










The composition is a function. 


This is a good reminder that a formula is not a function. There are always domain considerations.  In the case of a composition, we need to restrict the domain of the inner function to avoid inner function values that are not in the domain of the outer function.










\end{example}










\begin{example} Composition



\begin{itemize}
\item Let $g(x) = \frac{x-3}{x+1}$ with its natural or implied domain: $(-\infty, -1) \cup (-1. \infty)$. 

\item Let $H(t) = \frac{t+3}{1-t}$ with its natural or implied domain: $(-\infty, 1) \cup (1. \infty)$.
\end{itemize}



\[
(g \circ H)(y) = g(H(y)) = \frac{\left( \frac{y+3}{1-y} \right) - 3}{\left(  \frac{y+3}{1-y}\right) + 1} = \frac{y+3-3(1-y)}{y+3+(1-y)} = \frac{4y}{4} = y
\]



$g$ and $H$ are inverse functions.




\end{example}

Of course, that isn't the whole story.  As we said earlier, formulas are not functions. 

$(g \circ H)(r) = r$ and so $g \circ H = Id$, the identity function.  But remember our caveat.  There are always domain issues.  In this case $(g \circ H)(r) = r$ for almost all of the real numbers.

We know that, 


\[
(g \circ H)(r) = \frac{\left( \frac{r+3}{1-r} \right) - 3}{\left(  \frac{r+3}{1-r}\right) + 1} 
\]


This is equivalent to $r$ for almost all real numbers.


In this case, $(g \circ H)(r) = g(H(r))$, which means that $r \ne 1$, since $1$ is not in the domain of $H$.  Secondly, $-1$ is not in the domain of $g$.  We also cannot have values of $r$ that make $H(r) = -1$.




\[
H(r) = \frac{r+3}{1-r} = -1
\]

\[
r + 3 = -(1-r) = -1 + r
\]


\[
3 = -1
\]


There are no such real numbers.


Wait.  There is more to this story. We would like the other way as well: $(H \circ g)(z) = z$ and so $H \circ g = Id$.


Thus, $z \ne -1$.  We must also avoid $g(z) = 1$, because $1$ is not in the domain of $H$.



\[
g(z) = \frac{z-3}{z+1} = 1
\]

\[
z - 3 = z + 1 
\]


\[
-3 = 1
\]


There are no such real numbers.













$\blacktriangleright$ Therefore, $g$ and $H$ are inverse functions on $(-\infty, 1) \cup (-1, 1) \cup (1,\infty)$.




















\section*{The Other Half of the Story}


We would like our function arithmetic to mimic our arithmetic for numbers.  For numbers, the inverses are commutative.

\[
4 + (-4) = (-4) + 4 = 0
\]


\[
4 \cdot \frac{1}{4} = \frac{1}{4} \cdot 4 = 1
\]



In our example, we also want $(H \circ g)(z) = H(g(z)) = z$. It does.


\[
(H \circ g)(z) = \frac{\left( \frac{z-3}{z+1} \right) + 3}{1 - \left(  \frac{z-3}{z+1}\right)} = \frac{z-3 + 3(z+1)}{z+1-(z-3)} = \frac{4z}{4} = z
\]


In this case, we cannot have $z = -1$






$\blacktriangleright$  Therefore, $g$ and $H$ are inverse functions on $(-\infty, -1) \cup (-1, 1) \cup (1, \infty)$.























\begin{center}
\textbf{\textcolor{green!50!black}{ooooo-=-=-=-ooOoo-=-=-=-ooooo}} \\

more examples can be found by following this link\\ \link[More Examples of Composition]{https://ximera.osu.edu/csccmathematics/precalculus/precalculus/composition/examples/exampleList}

\end{center}









\end{document}
